\chapter{اصلاحات رسانه‌ای و نظام آموزشی}
\label{app:media_edu}

پایداری دموکراسی نیازمند شهروندانی آگاه و رسانه‌هایی مستقل است.

\section{اصلاح نظام آموزشی: اولویت استراتژیک}

نظام آموزشی باید از یک ابزار ایدئولوژیک به بستری برای تربیت شهروندان منتقد و خلاق تبدیل شود.

\begin{modernbox}{اقدامات فوری آموزشی (ماه اول)}
\begin{itemize}
    \item حذف محتوای ایدئولوژیک و سیاسی از کتاب‌های درسی.
    \item آزادی پوشش در فضاهای آموزشی و لغو جداسازی جنسیتی اجباری.
    \item استقلال کامل دانشگاه‌ها و لغو نظارت‌های عقیدتی بر اساتید و دانشجویان.
    \item آموزش زبان مادری در مناطق دو زبانه.
\end{itemize}
\end{modernbox}

\section{مدیریت رسانه در دوران گذار}

\begin{center}
\begin{tikzpicture}[node distance=2.5cm]
    \tikzstyle{block} = [rectangle, draw, fill=MainBlue!10, text width=6.5cm, text centered, rounded corners, minimum height=3.5em]
    \tikzstyle{line} = [draw, -Latex, line width=1pt]

    \node [block] (tv) {\textbf{۱. صداوسیمای ملی} \\ تبدیل به رسانه‌ی عمومی (Public Broadcaster) تحت نظارت پارلمانی};
    \node [block, below of=tv] (press) {\textbf{۲. مطبوعات آزاد} \\ تسهیل صدور مجوز و لغو سانسور پیشینی};
    \node [block, below of=press] (digital) {\textbf{۳. آزادی اینترنت} \\ رفع فوری فیلترینگ و تضمین حقوق دیجیتال};

    \path [line] (tv) -- (press);
    \path [line] (press) -- (digital);
\end{tikzpicture}
\end{center}

\section{چارت زمانی اصلاحات فرهنگی-آموزشی}

\begin{center}
\captionof{table}{مراحل اصلاحات فرهنگی}
\begin{tabular}{R{4cm} R{5cm} R{6cm}}
\hline
\textbf{بازه زمانی} & \textbf{اولویت آموزشی} & \textbf{اولویت رسانه‌ای} \\
\hline
%\endfirsthead
%\hline
%\textbf{بازه زمانی} & \textbf{اولویت آموزشی} & \textbf{اولویت رسانه‌ای} \\
%\hline
%\endhead
%\hline
%\endfoot
هفته اول & لغو حجاب اجباری و تفکیک جنسیتی & رفع فیلترینگ تلگرام، یوتیوب، اینستاگرام \\
ماه اول تا سوم & بازنویسی دروس تاریخ و بینش & بازگشایی انجمن صنفی روزنامه‌نگاران \\
سال اول & اصلاح ساختار عالی آموزش و پژوهش & تصویب قانون جامع آزادی اطلاعات \\
\hline
\end{tabular}
\end{center}
